%% LyX 2.3.4.2 created this file.  For more info, see http://www.lyx.org/.
%% Do not edit unless you really know what you are doing.
\documentclass[11pt,spanish]{article}
\renewcommand{\rmdefault}{cmr}
\usepackage[T1]{fontenc}
\usepackage[utf8]{inputenc}
\usepackage{geometry}
\geometry{verbose,tmargin=2cm,bmargin=2cm,lmargin=2cm,rmargin=2cm}
\setlength{\parskip}{\bigskipamount}
\setlength{\parindent}{0pt}
\usepackage{amstext}
\usepackage{amsthm}
\usepackage{amssymb}

\makeatletter
%%%%%%%%%%%%%%%%%%%%%%%%%%%%%% Textclass specific LaTeX commands.
\theoremstyle{definition}
\newtheorem{xca}{\protect\exercisename}[section]
\theoremstyle{definition}
\newtheorem{defn}{\protect\definitionname}[section]
\theoremstyle{definition}
\newtheorem{example}{\protect\examplename}[section]

%%%%%%%%%%%%%%%%%%%%%%%%%%%%%% User specified LaTeX commands.
\usepackage{titling}
\setlength{\droptitle}{-2cm}
\usepackage{multicol}
\usepackage{enumitem}
\usepackage{proof}
\usepackage{mathrsfs}
\usepackage{amsfonts}
\usepackage{forest}
\usepackage{ebproof}
\usepackage{amssymb}
\definecolor{Guinda}{rgb}{0.49, 0.05, 0.12}
\definecolor{Rojo}{rgb}{0.8, 0, 0}
\definecolor{Verde}{rgb}{0,0.4,0}
\definecolor{Azul}{rgb}{0,0.4,0.4}
\definecolor{Naranja}{rgb}{0.8,0.4,0}
\usepackage{tikz}
\usetikzlibrary{automata, positioning, arrows}

\makeatother

\usepackage{babel}
\addto\shorthandsspanish{\spanishdeactivate{~<>.}}

\usepackage{listings}
\lstset{mathescape=true}
\providecommand{\definitionname}{Definición}
\providecommand{\examplename}{Ejemplo}
\providecommand{\exercisename}{Ejercicio}
\renewcommand{\lstlistingname}{\inputencoding{latin9}Listado de c�digo}

\begin{document}
\title{Teoría de la Computación, 2021-1\\
Nota de clase 2: Gramáticas Regulares}
\author{Manuel Soto Romero}
\date{\today \\ Colegio de Ciencia y Tecnología UACM}

\maketitle
\renewcommand*{\thefootnote}{\arabic{footnote}}
\setcounter{footnote}{0}
\setcounter{section}{2}
\renewcommand*\lstlistingname{C\'odigo}

\subsection{Introducción a las gramáticas y lenguajes formales}

Como concluimos en la sección anterior, una vez que definimos un lenguaje
formal, nuestra mayor preocupación será si una cadena pertenece o
no a este. Dicho de otra forma, queremos ver si una cadena está bien
escrita de acuerdo a dicho lenguaje. Un ejemplo de lenguaje formal
que seguramente conociste en el curso de Matemáticas Discretas es
el Lenguaje de la Lógica Proposicional, mejor conocido como \textsc{LProp}.
Por cierto, es común que a las cadenas también se les llamen \emph{expresiones}.

El alfabeto de la Lógica proposicional es el siguiente:
\begin{itemize}
\item \emph{Constantes lógicas: }$\bot,\top$ (Otra notación es la de usar
$0$ y $1$ o también $F$ y $T$)
\item \emph{Variables proposicionales: $p,q,r,s,...,p_{1},q_{2},...$}
\item \emph{Conectivos lógicos: $\lnot,\lor,\land,\rightarrow,\leftrightarrow$}
\item \emph{Paréntesis}: $(,)$
\end{itemize}
Con este alfabeto, el Lenguaje de la Lógica Proposicional (\textsc{LProp})
se define de la siguiente manera:
\begin{enumerate}
\item Las constantes lógicas son expresiones de \textsc{LProp}.
\item Las variables proposicionales son expresiones de \textsc{LProp}.
\item Si $\varphi$ es una expresión de \textsc{LProp}, entonces $\left(\lnot\varphi\right)$
lo es también (\emph{negación}).
\item Si $\varphi$ y $\psi$ son expresiones de \textsc{LProp}, entonces
las siguientes también lo son (\emph{operadores lógicos}):
\begin{enumerate}
\item $\left(\varphi\lor\psi\right)$ \emph{disyunción}
\item $\left(\varphi\land\psi\right)$ \emph{conjunción}
\item $\left(\varphi\rightarrow\psi\right)$ \emph{condicional}
\item $\left(\varphi\leftrightarrow\psi\right)$ \emph{bicondicional}
\end{enumerate}
\item Son todas.
\end{enumerate}
A partir de este Lenguaje Formal, veamos si las siguientes cadenas
(expresiones) forman parte de él.
\begin{enumerate}
\item $\left(\lnot p\right)$
\begin{itemize}
\item Por la regla 2, $p$ está \textsc{LProp} puesto que es una variable
proposicional.
\item Por la regla 3 y dado que $p$ está \textsc{LProp}, entonces $\left(\lnot p\right)$
está \textsc{LProp}.
\end{itemize}
\item $\left(q\land r\right)$
\begin{itemize}
\item Por la regla 2, tanto $q$ como $r$ están en \textsc{LProp}
\item Por la regla 4b, y dado que $q$ y $r$ están en \textsc{LProp},
entonces $\left(q\land r\right)$ está en \textsc{LProp}.
\end{itemize}
\item $\left(\left(\lnot p\right)\rightarrow\left(q\land r\right)\right)$
\begin{enumerate}
\item Como vimos anteriormente, $p,q$ y $r$ están en \textsc{LProp}.
\item De la misma forma $\left(\lnot p\right)$ está en \textsc{LProp}.
\item Igualmente vimos que $\left(q\land r\right)$ está en \textsc{LProp}.
\item Por la regla 4c y porque tanto $\left(\lnot p\right)$ y $\left(q\land r\right)$
están en \textsc{LProp}, entonces $\left(\left(\lnot p\right)\rightarrow\left(q\land r\right)\right)$
está en \textsc{LProp}.
\end{enumerate}
\item $(\left((p\land q)\land r\right)\rightarrow\left(\lnot p\right))$
\begin{enumerate}
\item Como vimos anteriormente, $p,q$ y $r$ están en \textsc{LProp}.
\item Por la regla 4b y porque $p$ y $q$ están en \textsc{LProp}, entonces
$\left(p\land q\right)$ está en \textsc{LProp}.
\item Por la regla 4b y porque $\left(p\land q\right)$ y $r$ están en
\textsc{LProp}, entonces $\left(\left(p\land q\right)\land r\right)$
está en \textsc{LProp}.
\item Como vimos anteriormente $\left(\lnot p\right)$ está en \textsc{LProp}.
\item Por la regla 4c y porque $\left(\left(p\land q\right)\land r\right)$
y $\left(\lnot p\right)$ están en \textsc{LProp}, entonces $(\left((p\land q)\land r\right)\rightarrow\left(\lnot p\right))$
está en \textsc{LProp}.
\end{enumerate}
\end{enumerate}
\bigskip{}

\begin{xca}
Se deja como ejercicio al lector la verificación de que la fórmula
$\left(\lnot\left(p\land\left(q\lor r\right)\right)\right)$ siguiente
está en \textsc{LProp}:
\end{xca}
%

\subsubsection{Gramáticas}

Otra forma de mostrar cómo se construyen las expresiones es mediante
lo que se conoce como \emph{gramática formal}, que es un mecanismo
sencillo de especificación de reglas de construcción, llamadas \emph{producciones}
o \emph{reglas de reescritura, }con las cuales se pueden generar expresiones
de un lenguaje. 

Las producciones tienen la siguiente forma:

\[
\texttt{símbolo ::= cadena}
\]

El símbolo \texttt{::=} se lee como ``se reescribe como'' y al
aplicar una regla particular sustituimos el lado izquierdo de este
símbolo por la cadena que se encuentra del lado derecho. Usando este
notación podemos describir el lenguaje de la lógica proposicional
con la siguiente gramática: 

\bigskip{}

\inputencoding{latin9}\begin{lstlisting}[mathescape=true,xleftmargin=5cm]
S ::= E
E ::= var
E ::= const
E ::= $\vartriangleright$ E
E ::= E $\star$ E
E ::= (E)
var ::= p | q | ...
const ::= $\bot$ | $\top$
$\vartriangleright$ ::= $\lnot$
$\star$ ::= $\lor$ | $\land$ | $\rightarrow$ | $\leftrightarrow$
\end{lstlisting}
\inputencoding{utf8}
A una colección de reglas de reescritura como la anterior, le llamamos
\emph{gramática} porque nos describe la \emph{forma }o \emph{reglas
sintácticas} que deben tener las expresiones lógicas bien construidas.
Para producir expresiones (cadenas) correctas de acuerdo a la gramática,
empezamos con el símbolo de la izquierda del \texttt{::=} de la primera
regla, y utilizamos las reglas de reescritura, que nos dicen que en
cualquier momento podemos sustituir parte de (o toda) la expresión,
si en ella localizamos una subexpresión que corresponda al lado izquierdo
de cualquiera de las reglas y la sustituimos por el lado derecho.
Veamos algunos ejemplos.
\begin{itemize}
\item Proceso de generación de la expresión $p$
\begin{itemize}
\item $\texttt{S}$
\item $\texttt{E}$
\item $\texttt{var}$
\item $p$
\end{itemize}
\item Proceso de generación de la expresión $\left(\top\lor\bot\right)$
\begin{itemize}
\item $\texttt{S}$
\item $\texttt{E}$
\item $\left(\texttt{E}\right)$
\item $\left(\texttt{E}\star\texttt{E}\right)$
\item $\left(\texttt{const}\star\texttt{E}\right)$
\item $\left(\texttt{const}\lor\texttt{E}\right)$
\item $\left(\texttt{const}\lor\texttt{const}\right)$
\item $\left(\top\lor\texttt{const}\right)$
\item $\left(\top\lor\bot\right)$
\end{itemize}
\item Proceso de generación de la expresión $\lnot\left(p\land\left(q\lor r\right)\right)$
\begin{itemize}
\item $\texttt{S}$
\item $\texttt{E}$
\item $\vartriangleright\texttt{E}$
\item $\lnot\texttt{E}$
\item $\lnot\left(\texttt{E}\right)$
\item $\lnot\left(\texttt{E}\star\texttt{E}\right)$
\item $\lnot\left(\texttt{var}\star\texttt{E}\right)$
\item $\lnot\left(\texttt{var}\land E\right)$
\item $\lnot\left(\texttt{var\ensuremath{\land\left(\texttt{E}\right)}}\right)$
\item $\lnot\left(p\land\left(\texttt{E}\right)\right)$
\item $\lnot\left(p\land\left(\texttt{E}\star\texttt{E}\right)\right)$
\item $\lnot\left(p\land\left(\texttt{var}\star\texttt{E}\right)\right)$
\item $\lnot\left(p\land\left(\texttt{var}\lor\texttt{E}\right)\right)$
\item $\lnot\left(p\land\left(\texttt{var}\lor\texttt{var}\right)\right)$
\item $\lnot\left(p\land\left(q\lor\texttt{var}\right)\right)$
\item $\lnot\left(p\land\left(q\lor r\right)\right)$
\end{itemize}
\end{itemize}
Una secuencia de aplicación de reglas de reescritura, como las de
los ejemplos anteriores, se conocen como \emph{derivación} de una
expresión (cadena); esta expresión es la que figura en el último de
los pasos. Estas derivaciones pueden presentarse así o gráficamente
utilizando un \emph{árbol}.

\bigskip{}

\begin{xca}
Dada la siguiente gramática que define paréntesis bien balanceados:\bigskip{}
\end{xca}
\inputencoding{latin9}\begin{lstlisting}[mathescape=true,xleftmargin=5cm]
E ::= ()
E ::= (E)
E ::= EE
\end{lstlisting}
\inputencoding{utf8}
Realiza la derivación de las siguientes expresiones:
\begin{itemize}
\item $\texttt{(())}$
\item $\texttt{(()())}$
\end{itemize}

\subsubsection{Árboles de derivación}

Para representar arboles de derivación, consideramos lo siguiente:
\begin{itemize}
\item Cada nodo tiene un símbolo asociado.
\item Para que de un nodo salgan aristas se requiere que el símbolo que
está en ese nodo aparezca del lado izquierdo de alguna producción.
\item Las aristas apuntan a cada uno de los símbolos que aparecen del lado
derecho de la producción utilizada.
\item Es importante observar que un nodo tiene un único símbolo asociado.
\item Si el símbolo en un nodo se reescribe en una sucesión de tres símbolos,
deberán salir tres líneas de él, una por cada símbolo.
\item Aquellos nodos de los que no salen líneas son llamados \emph{hojas}.
\item Para determinar cuál es la expresión (cadena) que corresponde a un
determinado árbol, vamos listando las hojas del árbol de izquierda
a derecha. A la cadena que se obtiene de esta manera le vamos a llamar
el \emph{resultado del árbol}.
\end{itemize}
Veamos algunos ejemplos.
\begin{itemize}
\item $p$\begin{center}
\begin{forest}	
	[\texttt{S}
		[\texttt{E}
			[\texttt{var}
				[$p$]
			]
		]
	]
\end{forest}
\end{center}
\item $\left(\top\lor\bot\right)$\begin{center}
\begin{forest}	
	[\texttt{S}
		[\texttt{E}
			[\texttt{(}]
			[\texttt{E}
				[\texttt{E}
					[\texttt{const}
						[$\top$]
					]
				]
				[$\star$
					[$\lor$]
				]
				[\texttt{E}
					[\texttt{const}
						[$\bot$]
					]
				]
			]
			[\texttt{)}]
		]
	]
\end{forest}
\end{center}
\item $\lnot\left(p\land\left(q\lor r\right)\right)$\begin{center}
\begin{forest}	
	[\texttt{S}
		[\texttt{E}
			[$\vartriangleright$
				[$\lnot$]
			]
			[\texttt{E}
				[\texttt{(}]
				[\texttt{E}
					[\texttt{E}
						[\texttt{var}
							[$p$]
						]
					]
					[$\star$
						[$\land$]
					]
					[\texttt{E}
						[\texttt{(}]
						[\texttt{E}
							[\texttt{E}
								[\texttt{var}
									[$q$]
								]
							]
							[$\star$
								[$\lor$]
							]
							[\texttt{E}
								[\texttt{var}
									[$r$]
								]
							]
						]
						[\texttt{)}]
					]
				]
				[\texttt{)}]
			]
		]
	]
\end{forest}
\end{center}
\end{itemize}
\bigskip{}

\begin{xca}
Realiza la derivación de las expresiones del Ejercicio 1.5.
\end{xca}

\subsubsection{La Jerarquía de Chomksy}

La jerarquía de Chomsky es una clasificación jerárquica de distintos
tipos de gramáticas que generan lenguajes formales. Esta jerarquía
fue descrita por Noam Chomsky en 1956. La Jerarquía de Chomsky consta
de cuatro niveles:
\begin{itemize}
\item \textbf{Gramáticas de tipo 0} (sin restricciones), que incluye a todas
las gramáticas formales. 
\item \textbf{Gramáticas de tipo 1} (gramáticas sensibles al contexto) generan
los lenguajes sensibles al contexto. 
\item \textbf{Gramáticas de tipo 2} (gramáticas libres del contexto) generan
los lenguajes independientes del contexto. 
\item \textbf{Gramáticas de tipo 3} (gramáticas regulares) generan los lenguajes
regulares. 
\end{itemize}
A lo largo del curso estudiaremos cada uno de estos tipos de gramática,
de momento basta con que sepas acerca de su existencia.

\subsection{Lenguajes Regulares}

Los lenguajes regulares se encuentran en la clase más pequeña que
incluye a los lenguajes más simples. Por ejemplo, el conjunto de todos
los números binarios. De hecho deben su nombre a la \emph{regularidad}
o repeticiones que ocurren entre sus cadenas. Por ejemplo:

\[
L_{1}=\left\{ abc,cc,abab,abccc,ababc\right\} 
\]

La regularidad consiste en que sus palabras inician con repeticiones
de $ab$ seguidas de repeticiones de $c$. De hecho, los lenguajes
finitos son también lenguajes regulares por definición:

\[
L_{2}=\left\{ el,coche,verde\right\} 
\]

La combinación de lenguajes regulares (unión o concatenación) también
producen un lenguaje regular. En general, diremos que un lenguaje
$L$ es regular si y sólo si se cumple al menos una de las siguientes
condiciones:
\begin{itemize}
\item $L$ es finito
\item $L$ es la unión o la concatenación de otros lenguajes regulares $R_{1}$
y $R_{2}$, $L=R_{1}\cup R_{2}$ o $L=R_{1}R_{2}$ respectivamente.
\item $L$ es la cerradura de Kleene de algún lenguaje regular, $L=R^{*}$.
\end{itemize}
Por ejemplo, sea el lenguaje $L$ de palabras formadas por $a$ y
$b$ que empiezan con $a$, tales como $aab,ab,a,abaa,$etc. Podemos
probar que este lenguaje es regular como sigue:

El alfabeto es $\Sigma=\left\{ a,b\right\} $. El lenguaje $L$ puede
ser visto como la concatenación de una $a$ con cadenas cualesquiera
de $a$ y $b$; ahora bien, éstas últimas son elementos de $\left\{ a,b\right\} ^{*}$,
mientras que el lenguaje que sólo contiene la palabra $a$ es $\left\{ a\right\} $.
Ambos lenguajes son regulares, es decir, $\left\{ a\right\} $ es
finito, por lo tanto regular, mientras que $\left\{ a,b\right\} ^{*}$
es la cerradura de $\left\{ a,b\right\} $, que es regulare por ser
finito también. Entonces la concatenación es $\left\{ a\right\} \left\{ a,b\right\} ^{*}$,
es regular.

En esta nota estudiaremos a los lenguajes regulares mediante algunos
mecanismos como son las gramáticas, los autómatas finitos y las expresiones
regulares.

\subsection{Gramáticas Regulares}

Aunque en la nota pasada usamos la notación $S:=E$ para representar
las reglas de producción, ésta es más utilizada en el ámbito del diseño
de lenguajes de programación. En Teoría de la Computación se prefiere
la notación $S\rightarrow E$. 

Decimos entonces que las gramáticas regulares son gramáticas cuyas
reglas son de la forma $A\rightarrow aB$ o bien $A\rightarrow a$,
donde $A$ y $B$ son símbolos no terminales y $a$ es un símbolo
terminal. Por ejemplo:

\[
\begin{array}{rcl}
S & \rightarrow & aA\\
S & \rightarrow & bA\\
A & \rightarrow & aB\\
A & \rightarrow & bB\\
A & \rightarrow & a\\
B & \rightarrow & aA\\
B & \rightarrow & bA
\end{array}
\]

En este caso, dado que en los símbolos no terminales se encuentran
a la derecha, decimos que es una \emph{gramática} \emph{regular derecha},
y cuando colocamos el símbolo no terminal a la izquierda, diremos
evidentemente que, es \emph{regular izquierda. }Por ejemplo:

\[
\begin{array}{rcl}
S & \rightarrow & Aa\\
S & \rightarrow & Ab\\
A & \rightarrow & Ba\\
A & \rightarrow & Bb\\
A & \rightarrow & a\\
B & \rightarrow & Aa\\
B & \rightarrow & Ab
\end{array}
\]

\bigskip{}

\begin{defn}
\emph{(Gramática Regular) Una gramática regular es una cuadrupla $\left\langle V,\Sigma,R,S\right\rangle $
en donde:}
\end{defn}
\begin{itemize}
\item $V$ \emph{es el alfabeto de variables}
\item $\Sigma$ \emph{es el alfabeto de constantes}
\item $R$ \emph{es el conjunto de reglas de producción que además es un
subconjunto de $V\times\left(\Sigma V\cup\Sigma\right)$}
\item $S$ \emph{es el símbolo inicial, un elemento de $V$}
\end{itemize}
%
\textbf{Derivación}
\begin{itemize}
\item Una cadena $uXu$ deriva en un paso una cadena $u\alpha v$, escrito
como $uXv\Rightarrow u\alpha v$, si hay una regla de producción $X\rightarrow\alpha\in R$
en la gramática. 
\item Una cadena $w\in\Sigma^{*}$ (formada sólo por símbolos terminales)
es derivable a partir de una gramática $G$ si existe una secuencia
de pasos de derivación $S\Rightarrow\alpha_{1}\Rightarrow...\Rightarrow w$.
\item A una secuencia de pasos de derivación le llamamos simplemente derivación.
\item Dicho de otra forma, una palabra $w\in\Sigma^{*}$ es derivable a
partir de $G$ si podemos llegar a ella a partir de $S$ y escribimos

\[
S\Rightarrow^{*}w
\]

\end{itemize}
\bigskip{}

\begin{defn}
\emph{(Lenguaje de una Gramática) El Lenguaje generado por una gramática
$G$, $L\left(G\right)$, es igual al conjunto de las cadenas derivables
a partir de su símbolo inicial. Es decir:}

\[
L\left(G\right)=\left\{ w\in\Sigma^{\text{*}}:S\Rightarrow^{\text{*}}w\right\} 
\]
\end{defn}
%
Por ejemplo, el lenguaje reconocido por nuestra primera gramática
es el de las palabras en $\left\{ a,b\right\} $ de longitud par terminadas
en $a$. Veamos otro ejemplo.

\bigskip{}

\begin{example}
Diseñar una gramática que genere el lenguaje de las palabras en $\left\{ a,b\right\} $
que contiene la subcadena $bb$, tales como $abb$, $ababba$, etc.

\emph{Solución:}

Podemos usar los símbolos no terminales para \emph{recordar} situaciones.
Es decir:
\end{example}
\begin{itemize}
\item Un símbolo $A$ que recuerda que aún no se ha producido ninguna $b$.
\item Un símbolo $B$ que recuerda que se produjo una $b$.
\item Un símbolo $C$ que recuerda que ya se produjeron las dos $b$.
\end{itemize}
\[
\begin{array}{rcl}
A & \rightarrow & aA\\
A & \rightarrow & bB\\
B & \rightarrow & aA\\
B & \rightarrow & bC\\
B & \rightarrow & b\\
C & \rightarrow & aC\\
C & \rightarrow & bC\\
C & \rightarrow & a\\
C & \rightarrow & b
\end{array}
\]

\bigskip{}

\begin{xca}
Diseña una gramática que dado el siguiente alfabeto $\left\{ a,b,c,1,2,3\right\} $
permita construir identificadores válidos. Decimos que un identificador
es válido si éste es una combinación de letras y números pero con
la restricción de que nunca debe empezar por un número. Por ejemplo
$abc$, $a12$, $a1b$.
\end{xca}
%
\begin{thebibliography}{1}
\bibitem{key-1}Favio E. Miranda, Elisa Viso, \emph{Matemáticas Discretas,
}Las Prensas de Ciencias, Facultad de Ciencias UNAM, 2010.
\end{thebibliography}

\end{document}
